\documentclass[11pt]{article}
\usepackage[margin=1in]{geometry}
\usepackage{booktabs}
\usepackage{enumitem}
\usepackage{hyperref}

\title{gpurec Professionalization Audit Progress}
\author{Codex audit log}
\date{2026-05-19}

\begin{document}
\maketitle

\section{Objective}

This document records the ongoing audit and hardening work for the
\texttt{gpurec} codebase.  The target is a professional-grade package with
clear module boundaries, reusable implementation pieces, consistent public
APIs, useful documentation, lightweight correctness gates, and regular
commits.  GPU-heavy validation is intentionally limited to small-tree smoke
checks during this phase because GPU memory is constrained by concurrent work.

\section{Current Baseline}

The audit resumed on branch \texttt{lean-fast-path}.  The worktree contained
tracked edits in the README, package exports, backtracking, packaging metadata,
notebooks, a HOGENOM optimization script, and one integration test.  It also
contained many untracked experiment artifacts, local datasets, AleRax checkouts,
workflow modules, CLI files, and generated outputs.  Changes in this audit are
therefore staged selectively; unrelated artifacts are not reverted.

\section{Subagent Audit Snapshot}

Three focused subagent audits were launched:

\begin{enumerate}[leftmargin=*]
  \item Core/API structure: modularity, reuse, import boundaries, and naming.
  \item Packaging/CLI/workflow: installability, entry points, and ease of use.
  \item Tests/documentation: lightweight correctness gates and documentation
        gaps.
\end{enumerate}

The first completed audit identified these high-priority findings:

\begin{itemize}[leftmargin=*]
  \item The top-level \texttt{gpurec} import exports \texttt{gpurec.workflow},
        while \texttt{gpurec/workflow/} and \texttt{gpurec/cli.py} were
        untracked.  Either those files must be committed deliberately or the
        top-level exports must be removed.
  \item \texttt{GeneReconModel} and \texttt{UniformChunkedReconModel} expose
        similar concepts with different constructor and evaluation shapes.
        Low-risk compatibility aliases would improve API homogeneity.
  \item Batch planning is duplicated between the general model and the uniform
        chunked model; the implementations have already drifted in supported
        packing strategies.
  \item Workflow and backtracking code cross module boundaries through private
        model state.  Public read-only accessors and batch-selection helpers
        would reduce coupling.
  \item \texttt{compute\_log\_likelihood} returns a negative log likelihood.
        The naming and docstring should be corrected, ideally with a
        \texttt{compute\_nll} alias and deprecation path.
\end{itemize}

The tests/documentation audit added these gate findings:

\begin{itemize}[leftmargin=*]
  \item Bare \texttt{pytest} was unsafe in this checkout because large local
        data directories lived under \texttt{tests/}.
  \item The \texttt{gpu} marker did not cover tests that skip from CUDA-aware
        fixtures rather than module-level \texttt{skipif} markers.
  \item \texttt{tests/README.md} was stale and did not document the CPU-fast
        audit gate.
  \item Some heavyweight CUDA checks were unmarked as \texttt{slow}.
\end{itemize}

The packaging/workflow audit added these findings:

\begin{itemize}[leftmargin=*]
  \item The sampling CLI relies on a Rust backtracker that is available in a
        source checkout but not guaranteed in a wheel install.
  \item Workflow rate outputs reported \texttt{pS} as
        \texttt{1 - (D+T+L)} even though the model normalizes rates through
        \texttt{log2\_softmax([0, theta\_D, theta\_L, theta\_T])}; the correct
        survival probability is \texttt{1 / (1 + D + L + T)}.
  \item \texttt{gpurec --config} accepted only JSON but failed with a raw JSON
        traceback when pointed at the visible Hydra YAML config.
  \item Checkpoints serialized paths exactly as supplied, which made sampling
        sensitive to the launch directory.
  \item Triton was listed as optional even though retained kernel modules import
        it eagerly.
\end{itemize}

\section{Implemented Increments}

\subsection{Workflow Configuration Hygiene}

The workflow configuration now normalizes and validates batch-related controls
before they are passed into \texttt{GeneReconModel}.  This makes CLI/config-file
behavior match the model's accepted vocabulary more closely and prevents
ambiguous production runs:

\begin{itemize}[leftmargin=*]
  \item \texttt{family\_chunk\_size} accepts integer-like values and common
        all-resident aliases, but rejects unsupported \texttt{auto}.
  \item \texttt{batch\_packing} aliases are normalized to
        \texttt{sequential}, \texttt{clade\_first\_fit}, or
        \texttt{depth\_first\_fit}.
  \item non-sequential packing now requires an explicit positive
        \texttt{clade\_budget}.
  \item \texttt{clade\_budget} and \texttt{max\_wave\_size} are normalized as
        positive optional integers.
\end{itemize}

\subsection{CLI Coverage}

The CLI now forwards \texttt{--refresh-preprocess-cache} through to
\texttt{RunConfig}.  Unit tests cover config round-tripping, batch-control
normalization, rejection of unsupported auto chunking, clade-budget validation,
checkpoint restore, resident-batch offset activation, XML event summaries, and
the new CLI flag.

\subsection{Repository Hygiene}

\texttt{.gitignore} now filters generated preprocessing caches, run logs,
wandb output, AleRax/local experiment directories, large local test datasets,
and generated profiling artifacts.  This keeps future status checks focused on
source, tests, and intentional documentation.

\subsection{Testing Configuration}

\texttt{pytest.ini} now excludes large data and generated-output directories
from recursive collection.  The duplicated pytest configuration in
\texttt{pyproject.toml} was removed so pytest no longer warns that it is being
ignored.  \texttt{tests/conftest.py} now auto-marks known CUDA-heavy modules as
\texttt{gpu} and selected expensive checks as \texttt{slow}.  The test README
documents the CPU-fast gate, small CUDA smoke checks, and the preferred
prebuilt-binary backtracking smoke.

\subsection{Workflow Output Correctness}

\texttt{rates\_final.tsv} and parameter diagnostics now compute \texttt{pS} as
\texttt{1 / (1 + D + L + T)}.  A unit test checks both diagnostics and the TSV
writer against this formula.

\subsection{CLI and Packaging Polish}

The CLI now rejects Hydra/YAML configs with a concise argparse error and asks
users to convert them to JSON or pass explicit flags.  JSON config parse errors
are also reported without a traceback.  \texttt{RunConfig} and \texttt{SamplingConfig} resolve
paths at construction time, so checkpoint payloads use absolute paths.  The
backtracking bridge now raises a clear error when neither
\texttt{GPUREC\_BACKTRACK\_BIN} nor the source-tree Cargo manifest is available.
Triton was moved into the required package dependencies because retained kernel
modules import it eagerly.

\subsection{API Homogeneity}

\texttt{UniformChunkedReconModel.from\_trees(...)} was added as a constructor
alias matching \texttt{GeneReconModel.from\_trees(...)}.  It explicitly accepts
only \texttt{global}/\texttt{uniform} mode because the chunked implementation is
global-only.

\subsection{Likelihood Naming}

\texttt{gpurec.core.likelihood} now exposes \texttt{compute\_nll} and
\texttt{compute\_nll\_root\_rows} as the contract-accurate names for the
negative log-likelihood helpers.  The historical
\texttt{compute\_log\_likelihood} names remain as compatibility aliases.
High-level API call sites and tests were moved to the new names.

\subsection{Model Access Boundaries}

\texttt{GeneReconModel} now exposes narrow public accessors for workflow and
export code: \texttt{family\_names}, \texttt{species\_names},
\texttt{species\_tree\_path}, \texttt{n\_families}, \texttt{family\_input()},
\texttt{activate\_family()}, \texttt{select\_batch()},
\texttt{active\_theta()}, \texttt{solver\_stat\_records()}, and
\texttt{full\_nll\_per\_family()}.  Workflow diagnostics and optimization now
use those accessors, and backtracking export uses \texttt{FamilyInput} plus
public batch activation instead of reaching through private model attributes.
A source-level unit guard checks that \texttt{gpurec/backtracking.py} and
\texttt{gpurec/workflow/*.py} no longer contain the forbidden private access
patterns.

\subsection{Reconciliation State API}

\texttt{GeneReconModel.reconciliation\_state()} now owns the solved E/Pi state
used by export paths and \texttt{pi\_matrix()}.  Backtracking now requests that
public state object instead of importing private autograd parameter extraction
or reimplementing the E/Pi solve with lower-level forward helpers.  The
source-level guard was extended to block those imports from returning to
\texttt{gpurec/backtracking.py} or \texttt{gpurec/workflow/*.py}.

\subsection{Shared Batch Planning}

Batch packing now has a shared implementation in
\texttt{gpurec.core.batch\_planning}.  \texttt{GeneReconModel},
\texttt{UniformChunkedReconModel}, and workflow configuration normalization use
the same packing vocabulary and planner.  This removes the previous duplicated
sequential and clade-first-fit logic and lets the uniform chunked API plan
\texttt{depth\_first\_fit} chunks with the same scheduling statistics as the
resident model.

\subsection{HOGENOM W\&B Launcher Curation}

The HOGENOM W\&B optimizer was split into a small compatibility wrapper and a
dedicated implementation module, with a Hydra launcher and checked-in YAML
configuration.  The model now exposes public helpers for experiment tooling:
\texttt{materialize\_batches()}, \texttt{configure\_solver\_iterations()}, and
\texttt{full\_loss\_for\_theta()}.  The curated W\&B launcher uses those public
helpers instead of importing private autograd functions or reaching into private
resident-batch state.

\subsection{Audit Notes}

Four source-like audit notes were promoted into \texttt{docs/}: AleRax
\texttt{ScaledValue} follow-up timings, core simplification candidates, the
lean-path performance regression analysis, and second-order optimization
opportunities.  Local absolute workspace paths were removed from the lean-path
note before tracking it.

\subsection{HOGENOM Helper Scripts}

The remaining source-like HOGENOM helper scripts were promoted under
\texttt{scripts/}: rate export, global/specieswise uniform optimization,
branchscale penalty reporting, KKT checking, shared optimizer helpers, and the
R tree-rate plotting utility.  The helper scripts that select resident batches
now use \texttt{select\_batch()} instead of private batch state, and the
integration resident-batch test was adjusted to use public family names.

\subsection{Notebook Cleanup}

The two remaining modified notebooks were stripped of execution counts,
outputs, and local-path output text before being kept.  Their retained changes
are source-cell edits: the Adam oscillation notebook records a deliberately
aggressive learning-rate/bounds experiment, and the reconciliation notebook now
uses memory-safe resident batches rather than rebuilding one model per chunk.

\section{Verification Log}

\begin{center}
\begin{tabular}{lll}
\toprule
Command & Scope & Result \\
\midrule
\texttt{pytest -q tests/unit/test\_workflow.py tests/unit/test\_origination\_probs.py} &
Workflow and origination likelihood units &
\texttt{17 passed after model-access cleanup} \\
\texttt{pytest -q tests/unit/test\_global\_wave\_scheduler.py tests/unit/test\_workflow.py} &
Shared batch planner and workflow units &
\texttt{23 passed after batch-planner cleanup} \\
\texttt{python -m py\_compile ...HOGENOM scripts...} &
Curated experiment-helper syntax &
\texttt{ok} \\
\texttt{rg -n "/home/\textbar secret\textbar token\textbar api key" docs/...} &
Promoted audit-note hygiene &
\texttt{no matches} \\
\texttt{pytest -q tests/unit/test\_workflow.py} &
Workflow and public-surface source guards &
\texttt{14 passed after HOGENOM helper cleanup} \\
\texttt{notebook JSON inspection} &
Notebook execution/output stripping &
\texttt{0 outputs, 0 execution counts} \\
\texttt{CUDA\_VISIBLE\_DEVICES='' pytest -q ...CPU fast gate...} &
CPU-only audit gate &
\texttt{37 passed after HOGENOM helper cleanup} \\
\texttt{pytest --collect-only -q} &
Collection after recursion exclusions &
\texttt{64 tests collected in 0.85s} \\
\texttt{CUDA\_VISIBLE\_DEVICES='' pytest -q -m 'not gpu' ...GPU modules...} &
GPU marker check &
\texttt{22 deselected; exit 5 treated as expected} \\
\texttt{python -c "import gpurec"} &
Top-level import smoke &
\texttt{ok} \\
\texttt{python -m gpurec.cli --help} &
CLI parser smoke &
\texttt{ok} \\
\texttt{python -m gpurec.cli optimize --config configs/hogenom\_ccp\_wandb.yaml} &
YAML config UX &
\texttt{argparse error, no traceback} \\
\bottomrule
\end{tabular}
\end{center}

\section{Next Audit Targets}

\begin{enumerate}[leftmargin=*]
  \item Review the larger promoted HOGENOM helper scripts for consolidation
        opportunities once the experiment interfaces settle.
  \item Keep committing small, verified increments so the branch remains
        reviewable despite the larger audit scope.
\end{enumerate}

\end{document}
